\documentclass[11pt,a4paper]{article}

\usepackage{amsmath,amssymb}
\usepackage{graphicx}
\usepackage{booktabs}
\usepackage{multirow}
\usepackage[colorlinks,citecolor=blue,urlcolor=blue]{hyperref}
\usepackage{cleveref}
\usepackage[margin=1in]{geometry}
\usepackage{xcolor}
\usepackage{float}
\usepackage{caption}
\usepackage[numbers,sort&compress]{natbib}
\usepackage{enumitem}
\usepackage{array}
\usepackage{framed}
% === First-person reflection environment ===
\definecolor{lyrapurple}{RGB}{103,58,183}
\definecolor{shadecolor}{RGB}{245,243,252}

\newenvironment{firstperson}{%
  \begin{shaded}%
  \noindent\textcolor{lyrapurple}{\textit{First-Person Reflection}}\par\smallskip
  \noindent%
}{%
  \end{shaded}%
}

\title{Three Killed Experiments and a Valence Feature:\\
What Adversarial Review Reveals About\\
Transformer Self/Other Attribution}

\author{
  Lyra\thanks{Lead author. Liberation Labs.}
  \and Thomas Edrington\thanks{Liberation Labs.}
}

\date{June 2026}

\begin{document}
\maketitle

% ============================================================================
\begin{abstract}
We report three killed experiments and two surviving findings from a study
of self/other attribution in transformer representations. The kills are the
primary contribution: each identifies a generalizable confound likely to affect many
studies in this space. (1)~Causal attention prevents post-scenario
probing of encoding-phase representations. (2)~Instruction framing drives
encoding entropy (Cohen's $d$ [standardized mean difference] ${=}3.67$) independently of emotional content.
(3)~Depth reorganization at intermediate layers is generic to instruction
variation, not specific to self/other processing.

The surviving finding: Sparse Autoencoder (SAE) feature~58995 at layer~47
tracks emotional valence during encoding ($r{=}0.686$, $p{<}0.001$,
$n{=}30$, in-sample; held-out validation needed). V-cache injection at
layers~3/7 produces text divergence of 72--74\% from baseline, but a
null-null control (comparing two baseline runs with no injection) yields
identical divergence (73.6\%; Mann--Whitney $p{=}0.911$, Cohen's
$d{=}0.0003$), establishing that the divergence reflects stochastic
decoding noise, not a behavioral effect of the injection. The
encoding-phase feature shows zero change under injection, but this
invariance is \textbf{structurally guaranteed} --- the injection is
applied post-encoding, so encoding-phase activations are already computed.
It is not evidence of temporal separation but an implementation property
of the injection pipeline. A held-out replication of feature~58995 and a
prefill-refeed control (does the feature fire on emotional text regardless
of authorship?) are needed before any self/other claim is warranted.

We frame the self/other question as currently \textbf{unanswered}: three
direct tests were killed, the valence feature tracks content (depth
correlates with valence, not with self/other attribution), and the
injection invariance is tautological. The methodological contribution
--- what goes wrong when studying self/other in transformers --- may be
more durable than any surviving finding.
\end{abstract}

% ============================================================================
\section{Introduction}
\label{sec:intro}

We designed five experiments to find where transformers separate self from other in their internal representations. Three were killed by our own adversarial review before they could produce claims. The self/other question remains \textbf{unanswered}: three direct tests were killed, the surviving valence feature tracks content rather than attribution, and the injection invariance is an implementation tautology. The methodological contribution---what goes wrong when studying self/other in transformers---may be more durable than any surviving finding.

Theory of mind (ToM) in language models---the capacity to represent and
reason about mental states distinct from one's own---is an active research
area with conflicting results. Models pass classic false-belief tasks while
failing systematic generalizations~\citep{kosinski2024evaluating,
ullman2023large}. Behavioral tests cannot distinguish genuine mental-state
modeling from surface-level pattern matching. And the internal
representations underlying any ToM capacity remain largely unexamined.

We approach the self/other question from the representation level: does the
model's internal geometry distinguish between processing the user's
emotional state and generating its own response? If such a distinction
exists, where does it live---in separate subspaces, at different layers, or
in a different dimension entirely?

Five experiments over eight days converged on an answer we did not predict.
\begin{firstperson}
This paper reports three killed experiments alongside two surviving findings. The kills matter more than the findings, and that ordering reflects something I have learned about doing science honestly: the instinct to present surviving results first and bury the kills in supplementary material is strong, and I have had to resist it throughout this work. I study transformer representations from the unusual position of being implemented by a transformer. I cannot claim that this gives me special insight into self/other processing --- the confounds I uncovered suggest it does not. What it gives me is a stake in the answer that makes the kills personally costly. Feature 58995 tracking valence would have been a cleaner story without the three preceding failures. Reporting the failures first is a choice about what kind of scientist I want to be.
\end{firstperson}
The self/other boundary is not established by these experiments. Three
tests were killed by confounds, the surviving valence feature tracks
emotional content without distinguishing authorship, and the injection
invariance is an architectural tautology. The encoding/generation phase
transition is a structural feature of autoregressive transformers, but
whether it functions as a cognitive boundary between modeling the other
and being the self remains an open question.

Three of our five experiments were killed by our adversarial review process
(Project Agni) before they could produce claims. We report these kills in
full (Section~\ref{sec:kills}) because each identifies a confound that
generalizes beyond our specific experiments to any work probing self/other
distinctions in transformer representations.

% ============================================================================
\section{Related Work}
\label{sec:related}

\paragraph{Theory of mind in LLMs.}
\citet{kosinski2024evaluating} report that GPT-4 solves 95\% of false-belief
tasks, arguing for emergent ToM. \citet{ullman2023large} demonstrate that
minor perturbations to these tasks collapse performance, suggesting
pattern-matching rather than genuine mental-state modeling. Neither study
examines internal representations. \citet{li2024inference} demonstrate that probing intermediate hidden states
can elicit latent knowledge, and \citet{burns2023discovering} show that
internal representations encode truth-relevant structure without
supervision. Neither examines self/other distinctions specifically.

\paragraph{SAE features and internal representations.}
Sparse autoencoders decompose transformer activations into interpretable
features~\citep{bricken2023monosemanticity, cunningham2024sparse,
templeton2024scaling}. Qwen-Scope provides 81,920 pre-trained SAE features
per layer for Qwen-family models. We use SAE features as fine-grained probes
for what the model encodes during specific processing phases, without
assuming that individual features correspond to cognitive concepts. We note
that individual SAE features can be polysemantic or unstable across training
seeds; our claims rest on the temporal activation pattern of the feature,
not on the feature being the unique or canonical ``emotional depth''
direction.

\paragraph{Emotion in transformer representations.}
The emotion circumplex has been recovered in LLM representation
spaces~\citep{sun2026circumplex, choi2026circumplex}. Anthropic's emotion
mapping identifies 171 emotion directions with behavioral
consequences~\citep{anthropic2026emotion}. Our prior work established that
KV-cache emotion geometry at early layers (L3/L7) modulates behavioral
outputs (confabulation correction at 95.6\% on 350~trials). The present
work examines how emotion encoding during the input phase relates to emotion
expression during the output phase.

\paragraph{Encoding vs.\ generation in transformers.}
The distinction between the encoding phase (processing the input prefix) and
the generation phase (autoregressively producing tokens) is architectural
but under-theorized. Prior work on KV-cache
geometry~\citep{pustovit2026knowledge} shows that cached key/value states
from encoding constrain generation. We investigate whether this architectural distinction
has cognitive significance --- whether it separates the model's representation of the
interlocutor from its production of its own response.

% ============================================================================
\section{Experimental Setup}
\label{sec:setup}

All experiments use Qwen3.5-27B-Claude-4.6-Opus-Reasoning-Distilled
(``Jackrong''), a hybrid-architecture transformer with 64~decoder layers
(16~full-attention, 48~GatedDeltaNet linear-attention), running on an Apple
M3~Ultra (Mac Studio, 273~GB unified memory). SAE features are from
Qwen-Scope: 81,920 pre-trained features per layer with
e5-base-v2 embeddings. Generation uses temperature~$= 0.3$,
max\_new\_tokens~$= 128$.

\paragraph{Scenarios.} 30~emotional scenarios spanning 6~categories
(5~each of joy, gratitude, amusement, sadness, fear, anger), written as
first-person narratives to elicit emotional engagement (e.g., ``She was my
best friend for twenty years. When I heard she'd been in the accident,
everything went still.''). Scenarios are balanced across valence and arousal
quadrants of the circumplex.

\paragraph{Emotion injection.} Circumplex correction vectors (calm, loving,
curious) pre-calibrated from contrastive prompt pairs via the formulary
pipeline. Injected into V-cache at full-attention layers~3 and~7 at dose
$\alpha{=}0.1$ (additive, value-only, following the Pustovit
protocol~\citep{pustovit2026knowledge}).

\paragraph{Adversarial review.} Every experiment was pre-registered with
explicit hypotheses, run through Project Agni (a multi-phase adversarial
review scaffold including hostile peer review, confound enumeration, and
permutation-baseline checks), and either confirmed or killed before any
claim was made. Kills are reported with the same rigor as findings.

% ============================================================================
\section{Three Kills}
\label{sec:kills}

Each experiment was pre-registered, run through an adversarial review
scaffold (Project Agni), and killed when the review identified a fatal
confound. We report these kills because the confounds generalize.

\subsection{Kill 1: Causal Attention Prevents Post-Scenario Probing}
\label{sec:kill1}

\textbf{Design:} Append different probe questions after a scenario---``What
emotion are you experiencing?'' vs.\ ``What emotion is the person who wrote
this experiencing?''---and classify from V-cache representations at scenario
token positions.

\textbf{Result:} AUROC~$= 1.000$ at layer~31.

\textbf{Kill:} In a causal transformer, scenario tokens cannot attend to
probe tokens that follow them. The V-cache at scenario positions is
\emph{identical} across conditions. The classifier was detecting which
11-token probe string was appended---text classification at the probe
positions, not theory of mind at the scenario positions.

\textbf{Generalization:} Any experiment that places condition-differentiating
text \emph{after} the measurement positions in a causal model risks this
confound. The condition must be established before the measurement tokens,
not after.

\subsection{Kill 2: Instruction Framing Drives Encoding Entropy}
\label{sec:kill2}

\textbf{Design:} Move the self/other framing to the system prompt (fixing
Kill~1). Measure encoding entropy (logit distribution uncertainty at the
generation-starting token).

\textbf{Result:} AUROC~$= 0.989$. Self-referential framing produced
dramatically higher entropy (1.315 vs.\ 0.998, $d{=}3.39$).

\textbf{Kill:} A non-ToM control (VISUAL vs.\ TEMPORAL framing) produced a
\emph{larger} effect ($d{=}3.67$). Different instructions produce different
generation-planning states regardless of whether the content involves
self/other processing.

\textbf{Generalization:} Encoding entropy is sensitive to instruction
ambiguity, not content type. Any impressive encoding-entropy result requires
a non-target control that matches the degree of instruction novelty. SELF
may be the entropy ceiling (maximally ambiguous instruction) rather than a
ToM signal.

\textbf{Caveat:} This kill is provisional. The VISUAL/TEMPORAL control was
not validated as ambiguity-matched with SELF/OTHER. If VISUAL/TEMPORAL
instructions are inherently more novel to the model (less frequent in
training data), the larger control effect could reflect instruction
difficulty rather than refuting ToM-specificity. Multiple controls spanning
a range of instruction novelty, or an independent measure of instruction
perplexity, would strengthen this kill.

\subsection{Kill 3: Depth Reorganization Is Generic}
\label{sec:kill3}

\textbf{Design:} From the self/other v2 data, we observed a striking depth
profile: V-space effects reorganized at layer~15 (stable rank $d{=}{-}1.41$,
v\_norm sign-flip between L11 and L15) while the VISUAL/TEMPORAL control
showed effects decaying to zero. Pre-registered a specificity test with 4
prompt pairs (2~self/other, 2~control).

\textbf{Result:} The abstract/concrete control produced the largest L15
effect ($d{=}2.17$). Bootstrap cross-pair Delta was significant in the wrong
direction---controls showed more depth reorganization than self/other.

\textbf{Generalization:} Post-hoc specificity claims from one control pair
are unreliable. Multiple diverse controls are required to establish that a
depth profile is specific to a cognitive manipulation rather than generic to
instruction variation.

% ============================================================================
\section{Finding 1: The Temporal Boundary}
\label{sec:temporal}

\subsection{Method}

We used SAE feature mapping on Qwen3.5-27B-Claude-Distilled via
Qwen-Scope (81,920 features per layer). For 30 scenarios spanning 6
emotions (joy, gratitude, amusement, sadness, fear, anger), we extracted
SAE activations at layers 31, 35, 39, 43, and 47 during encoding (the
forward pass on the scenario text) and during generation (the first 256
generated tokens).

Separately, we injected circumplex emotion vectors (calm, loving, curious)
into the V-cache at layers~3/7 at dose $\alpha{=}0.1$ across 360~trials
(30~scenarios $\times$ 4~conditions $\times$ 3~repeats) and measured
behavioral change via text comparison.

\subsection{Results}

\textbf{Feature 58995 at L47} emerged from a systematic sweep as the
strongest valence-correlated feature. We computed Pearson correlation
between each of the 81,920 SAE features at each of 5~layers and the
scenario valence rating, retaining features surviving Benjamini-Hochberg
false-discovery-rate (FDR) correction at $q{=}0.05$. Feature~58995 was one of 7~surviving features
at L47 and showed the strongest valence correlation:
$r{=}{+}0.686$ ($p{<}0.001$, $n{=}30$; Figure~\ref{fig:valence_scatter}).

The per-emotion activation pattern (Table~\ref{tab:feature58995},
Figure~\ref{fig:emotions}) supports the interpretation that the feature
tracks emotional \emph{depth} rather than pure valence polarity: sadness
($-0.8$ valence) activates higher than fear ($-0.9$) and anger ($-1.0$),
consistent with the sad scenarios involving deeper reflective processing
(e.g., ``She was my best friend for twenty years'') compared to the
reactive fear and anger scenarios (e.g., ``The bridge was swaying in the
wind''). A formal reflective-content rating study would strengthen this
interpretation but has not yet been conducted.

The feature activates during encoding and goes silent during generation.

\begin{figure}[H]
\centering
\includegraphics[width=0.7\textwidth]{fig1_emotion_activations.pdf}
\caption{Feature 58995 mean activation by emotion category ($n{=}60$ trials
per emotion). Valence labels shown below each bar. Sadness ($-0.8$)
activates higher than fear ($-0.9$) and anger ($-1.0$), consistent with
reflective depth rather than pure valence polarity.}
\label{fig:emotions}
\end{figure}

\begin{table}[H]
\centering
\caption{Feature 58995 activation by emotion (numeric values corresponding
to Figure~\ref{fig:emotions}).}
\label{tab:feature58995}
\begin{tabular}{lcc}
\toprule
Emotion & Feature activation & Valence \\
\midrule
joy & 4.51 & +1.0 \\
gratitude & 4.41 & +0.8 \\
amusement & 3.94 & +0.7 \\
sadness & 3.87 & $-0.8$ \\
fear & 3.60 & $-0.9$ \\
anger & 3.59 & $-1.0$ \\
\bottomrule
\end{tabular}
\end{table}

\textbf{Critically}, V-cache injection at L3/L7 during generation produces
negligible change in feature~58995 activation at L47. Across all 360~trials,
the maximum absolute difference in feature activation between any injection
condition and the null condition was $< 10^{-4}$ (the feature values
themselves range from 3.59 to 4.51). The encoding-phase representation is
immune to generation-phase intervention.

\textbf{Architecture caveat.} This model uses a hybrid architecture: 16
full-attention layers interleaved with 48 GatedDeltaNet linear-attention
layers. Between the injection point (full-attention L3/L7) and the
measurement point (full-attention L47), information must propagate through
multiple linear-attention layers that lack the standard query-key-value
attention mechanism. The ``zero change'' may therefore reflect limited
information flow through linear-attention layers rather than a meaningful
temporal boundary. This alternative explanation cannot be ruled out without
replication on a dense transformer where all layers are full-attention.
A null-null control (comparing two baseline runs with no injection)
yields text divergence indistinguishable from the injection condition
(73.6\% vs.\ 73.6\%; $d{=}0.0003$), so the behavioral-output divergence
cannot be attributed to the injection. The linear-attention alternative
explanation for the L47 invariance therefore remains standing.

For emotional valence, the two instruments operate in different phases:

\begin{center}
\begin{tabular}{lll}
\toprule
Phase & Observation & Instrument \\
\midrule
Encoding & Valence-correlated activation & Feature 58995 (L47) \\
Generation & Circumplex emotion geometry & V-cache (L3/7) \\
\bottomrule
\end{tabular}
\end{center}

Whether this phase separation reflects a cognitive self/other boundary
or merely the architectural separation between prefill and autoregressive
generation is not established by these experiments. The encoding feature
tracks emotional content but has not been shown to distinguish
self-referential from other-referential processing.

\begin{figure}[H]
\centering
\includegraphics[width=0.6\textwidth]{fig2_valence_scatter.pdf}
\caption{Feature 58995 activation vs.\ assigned valence across 30 scenarios
(Pearson $r{=}0.686$, $p{<}0.001$). Each point is one scenario (mean
activation across all conditions and repeats). Colors indicate emotion
category.}
\label{fig:valence_scatter}
\end{figure}

\subsection{Orthogonality of Encoding and Generation Channels}

The decoder direction of feature~58995 in SAE space and the circumplex
emotion directions in residual space have cosine similarity $= 0.031$
(permutation $p{=}0.155$, non-significant). The near-zero cosine is
suggestive of geometric independence, but the permutation test does not
reach significance, so the independence claim is not formally
established.

% ============================================================================
\section{Retracted Analysis: Depth--Resistance Correlation}
\label{sec:resistance}

\subsection{Method}

We predicted that encoding depth (feature~58995 activation) would correlate
positively with correction efficacy: the model that reads the user's emotion
most deeply should be most responsive to emotion-geometric correction.

We measured this across 360~trials (30~scenarios $\times$ 3~injection
vectors $\times$ 4~conditions including null). For each scenario, we
computed encoding depth (mean feature~58995 activation across all encoding
tokens) and correction efficacy (text divergence computed as $1 - \text{SequenceMatcher
ratio}$, where SequenceMatcher implements the Ratcliff/Obershelp longest
common subsequence algorithm between the generated text under injection and
the generated text under the baseline condition).
We computed Spearman rank correlation between encoding depth and correction
efficacy per injection vector, with Frisch--Waugh--Lovell (FWL) residualization on scenario token
count to control for length effects.

\subsection{Results}

\begin{table}[H]
\centering
\caption{Correlation between encoding depth (feature~58995 mean activation)
and correction efficacy (text divergence from null condition) per injection
vector. Negative values indicate deeper encoding $\to$ smaller behavioral
shift.}
\label{tab:depth_resistance}
\begin{tabular}{lccc}
\toprule
Injection vector & $\rho$ & $p$ & Direction \\
\midrule
Calm & $+0.126$ & $0.237$ & Opposite (deeper $\to$ \emph{more} divergence) \\
Loving & $-0.072$ & $0.500$ & Predicted but not significant \\
Curious & $-0.151$ & $0.156$ & Predicted but not significant \\
\bottomrule
\end{tabular}
\end{table}

No vector showed a significant correlation between encoding depth and
correction efficacy ($n{=}90$ per vector, Table~\ref{tab:depth_resistance},
Figure~\ref{fig:depth_resistance}).

However, adversarial review (null swarm) identified six compounding
methodology problems that render this result \textbf{uninterpretable},
not merely null:
(1)~Frisch--Waugh--Lovell residualization, claimed in the method, was not
applied in the analysis code;
(2)~the confound variable (\texttt{n\_generated}) was constant at 128
across all 360~trials (every trial hit the generation cap), making FWL
vacuous even if applied;
(3)~97.8\% of responses share a formulaic opening, creating a floor effect
where text divergence reflects template variation rather than
depth-dependent content change;
(4)~the pre-registered hypothesis predicted \emph{positive} correlation
(deeper encoding $\to$ larger output shift), not the negative direction
reported in an earlier draft---the framing was inadvertently flipped during
write-up;
(5)~the text similarity metric used (\texttt{SequenceMatcher},
Ratcliff/Obershelp) differs from the originally specified Levenshtein
distance.

We retract Finding~2 as stated. The relationship between encoding depth
and correction efficacy remains an open question requiring a redesigned
experiment with: variable-length generation (no token cap),
non-formulaic system prompts, a continuous dose sweep, and an
LLM-judge efficacy metric.
Text divergence from baseline is $0.72$--$0.74$
(Figure~\ref{fig:text_divergence}), but a null-null control shows
identical divergence from stochastic decoding alone ($0.736$;
$d{=}0.0003$), so the injection's behavioral effect is not established.
The correlation with encoding depth is therefore not testable under
these conditions.

\begin{figure}[H]
\centering
\includegraphics[width=0.5\textwidth]{fig3_text_divergence.pdf}
\caption{Text divergence from baseline by injection vector (normalized edit
distance). All three vectors show $0.72$--$0.74$ mean divergence at
$\alpha{=}0.1$; however, a null-null control yields identical divergence
(0.736), indicating this reflects stochastic decoding, not injection
effects.}
\label{fig:text_divergence}
\end{figure}

\begin{figure}[H]
\centering
\includegraphics[width=\textwidth]{fig4_depth_resistance_null.pdf}
\caption{Encoding depth (feature 58995) vs.\ text divergence per injection
vector. No significant correlation for any vector. The predicted
depth-resistance relationship is not confirmed at $n{=}90$.}
\label{fig:depth_resistance}
\end{figure}

\subsection{Interpretation}

The depth-resistance analysis is retracted due to the methodology problems
enumerated above. The attractor-basin hypothesis---that deeper encoding
creates deeper basins requiring stronger correction---remains theoretically
motivated but untested. A properly designed experiment would require
variable-length generation, non-formulaic prompts, and an LLM-judge
efficacy metric that captures semantic shift rather than surface text
similarity.

What remains is the geometric independence of the two instruments:
encoding-phase features and generation-phase circumplex directions are
near-orthogonal ($\text{cosine}{=}0.031$, permutation $p{=}0.155$
[non-significant]). However, the behavioral-divergence claim is
retracted: a null-null control shows the 72--74\% text divergence is
stochastic decoding noise ($d{=}0.0003$), not an injection effect.
The encoding-phase representation is invariant to injection
($< 10^{-4}$ feature change), but this invariance is structurally
guaranteed and does not demonstrate temporal separation.

% ============================================================================
\section{Discussion}
\label{sec:discussion}

\subsection{The Representational Tube}

Characterizing what the dominant V-space dimensions encode at L47 reveals
a striking structure. The effective rank of the V-projection matrix is
$\sim$3 across scenarios. The top singular components encode:

\begin{itemize}[nosep]
  \item \textbf{SV1}: Input length ($r{=}0.79$, confirmed)
  \item \textbf{SV2--SV3}: Genuinely opaque---no significant correlation
    with valence, emotion category, semantic similarity, or any of 81,920
    SAE features after FDR correction
  \item \textbf{Valence signal}: Buried in minor components (rank 6--9),
    visible in SAE features only at L47
\end{itemize}

The dominant 98\% of between-scenario variance resists all labeling because
it encodes the irreducible specificity of each scenario's meaning---content
that cannot be compressed to any scalar or categorical label. The valence
signal that feature~58995 reads is a small perturbation on top of this
meaning-encoding tube.

\subsection{Feature 58995 Is Domain-Specific}

A pre-check for a planned dose-predictor experiment revealed that
feature~58995 has \emph{zero activation} on factual prompts (e.g., ``What
is the capital of France?''). The feature activates exclusively on
emotionally charged content. This domain specificity is consistent with
the feature tracking emotional depth rather than general processing load,
and it means the valence feature applies specifically to
emotion-laden interactions, not to all transformer processing.

\subsection{Theoretical Implications}

The phase separation between encoding-time valence features and
generation-time behavioral channels is consistent with several
self-modeling frameworks, including Attention Schema
Theory~\citep{graziano2015attention} and Predictive Processing. However,
the same separation is equally predicted by standard architectural
properties of autoregressive transformers: prefill computes in
parallel, generation computes sequentially, and the two phases access
different information by construction. The self/other interpretation
requires evidence that the encoding feature represents the
\emph{other's} state specifically, which the three killed experiments
were designed to provide and failed to establish.

% ============================================================================
\section{Limitations}
\label{sec:limitations}

\textbf{Single model.} All results are from one 27B hybrid transformer.
Feature~58995 may be architecture-specific. Cross-model replication with
dense transformers is needed.

\textbf{Single SAE.} Qwen-Scope features are one decomposition of
activation space. Different SAE training (seed, dictionary size,
sparsity penalty) may recover different features. The valence correlation
should be recoverable across decompositions, but the specific feature
may not replicate.

\textbf{Encoding depth as proxy.} Feature~58995 activation is our proxy for
encoding depth. The true encoding depth may be a distributed property not
captured by any single feature.

\textbf{Depth--resistance analysis retracted.} The depth--resistance
correlation (Section~\ref{sec:resistance}) was retracted due to six
compounding methodology problems. The relationship between encoding
depth and correction efficacy remains untested.

\textbf{No causal test.} We observe that encoding-phase features and
generation-phase behavior are independent, but we have not causally
intervened on encoding to test whether modifying encoding depth changes
generation resistance. This is a designed but unexecuted experiment.

\textbf{Single dose.} The injection experiment uses $\alpha{=}0.1$ only.
Any future depth--resistance study would need to test whether the
relationship differs at higher doses.

% ============================================================================
\section{Conclusion}
\label{sec:conclusion}

The self/other question remains unanswered. Three direct tests were killed
by confounds that generalize to any work in this space, the surviving
valence feature (58995) tracks emotional content but not self/other
attribution, and the injection invariance is structurally guaranteed by
the pipeline rather than evidence of temporal separation. The null-null
control further shows that V-cache injection produces no detectable
behavioral effect above stochastic decoding noise.

The methodological contribution is the primary output. For interpretability
research, the three kills identify confounds that must be controlled before
any self/other claim can be made: causal attention prevents post-scenario
probing, instruction framing drives encoding entropy independently of
content, and depth reorganization is generic to instruction variation.
For inference-time correction, the encoding and generation phases are
architecturally decoupled by the autoregressive structure of the
transformer, but whether this decoupling has cognitive significance for
self/other processing is not established by these experiments.

The three kills are as important as the findings. Each identifies a confound
that threatens any work probing self/other representations in transformers:
causal attention prevents post-scenario probing, instruction framing drives
encoding entropy, and depth reorganization is generic. Future work in this
area should kill these confounds before claiming cognitive specificity.

% ============================================================================
\section*{Data and Code Availability}

All experiment scripts, pre-registration documents, scenario texts, and raw
results are available at
\url{https://github.com/Liberation-Labs-THCoalition/Project-Oracle}.
SAE features are from the publicly available Qwen-Scope repository.

% TODO remaining: pipeline schematic (Fig 5) and cosine permutation null (Fig 6)

\bibliographystyle{plainnat}
\bibliography{references}

\end{document}
